% Options for packages loaded elsewhere
\PassOptionsToPackage{unicode}{hyperref}
\PassOptionsToPackage{hyphens}{url}
\documentclass[
  ignorenonframetext,
]{beamer}
\newif\ifbibliography
\usepackage{pgfpages}
\setbeamertemplate{caption}[numbered]
\setbeamertemplate{caption label separator}{: }
\setbeamercolor{caption name}{fg=normal text.fg}
\beamertemplatenavigationsymbolsempty
% remove section numbering
\setbeamertemplate{part page}{
  \centering
  \begin{beamercolorbox}[sep=16pt,center]{part title}
    \usebeamerfont{part title}\insertpart\par
  \end{beamercolorbox}
}
\setbeamertemplate{section page}{
  \centering
  \begin{beamercolorbox}[sep=12pt,center]{section title}
    \usebeamerfont{section title}\insertsection\par
  \end{beamercolorbox}
}
\setbeamertemplate{subsection page}{
  \centering
  \begin{beamercolorbox}[sep=8pt,center]{subsection title}
    \usebeamerfont{subsection title}\insertsubsection\par
  \end{beamercolorbox}
}
% Prevent slide breaks in the middle of a paragraph
\widowpenalties 1 10000
\raggedbottom
\AtBeginPart{
  \frame{\partpage}
}
\AtBeginSection{
  \ifbibliography
  \else
    \frame{\sectionpage}
  \fi
}
\AtBeginSubsection{
  \frame{\subsectionpage}
}
\usepackage{iftex}
\ifPDFTeX
  \usepackage[T1]{fontenc}
  \usepackage[utf8]{inputenc}
  \usepackage{textcomp} % provide euro and other symbols
\else % if luatex or xetex
  \usepackage{unicode-math} % this also loads fontspec
  \defaultfontfeatures{Scale=MatchLowercase}
  \defaultfontfeatures[\rmfamily]{Ligatures=TeX,Scale=1}
\fi
\usepackage{lmodern}
\ifPDFTeX\else
  % xetex/luatex font selection
\fi
% Use upquote if available, for straight quotes in verbatim environments
\IfFileExists{upquote.sty}{\usepackage{upquote}}{}
\IfFileExists{microtype.sty}{% use microtype if available
  \usepackage[]{microtype}
  \UseMicrotypeSet[protrusion]{basicmath} % disable protrusion for tt fonts
}{}
\makeatletter
\@ifundefined{KOMAClassName}{% if non-KOMA class
  \IfFileExists{parskip.sty}{%
    \usepackage{parskip}
  }{% else
    \setlength{\parindent}{0pt}
    \setlength{\parskip}{6pt plus 2pt minus 1pt}}
}{% if KOMA class
  \KOMAoptions{parskip=half}}
\makeatother
\setlength{\emergencystretch}{3em} % prevent overfull lines
\providecommand{\tightlist}{%
  \setlength{\itemsep}{0pt}\setlength{\parskip}{0pt}}
\usepackage{bookmark}
\IfFileExists{xurl.sty}{\usepackage{xurl}}{} % add URL line breaks if available
\urlstyle{same}
\hypersetup{
  hidelinks,
  pdfcreator={LaTeX via pandoc}}

\author{\texorpdfstring{}{}}
\date{}

\begin{document}

\begin{frame}[fragile]{Abstract}
\protect\phantomsection\label{abstract}
This note connects ideal-gas and reference-liquid densities to a
reproducible, openly specified \textbf{composite mixture model} for
arthropod-scale effective density. Numbers reported here are
\textbf{model outputs} under stated assumptions: the mass-fraction
harmonic mixing rule developed in the ``Insect composite model''
section, tabulated component densities for cuticle, soft tissue,
hemolymph, and a phenomenological \textbf{internal gas} compartment, and
fixed handbook-style fluid anchors {[}@haynes2014crc{]}. They are not
species-specific measurements.

Under the built-in illustrative presets, central mixture densities fall
near \textbf{900--1030 kg m⁻³}, i.e.~within a few percent of fresh water
at 25 °C (997 kg m⁻³ in \texttt{fluid\_reference}). Ideal dry air at
standard temperature and pressure evaluates to about \textbf{1.29 kg
m⁻³} from \texttt{ideal\_gas\_density\_kg\_m3}, bracketed by a
documented literature band \textbf{1.18--1.32 kg m⁻³} for sanity
checks---not a fitted uncertainty on the insect model. Sensitivity to
the uncertain internal tracheal/air \textbf{mass} fraction is summarized
as a closed interval by rescaling non-gas compartments while sweeping
\texttt{internal\_gas}; this is a \textbf{parameter box}, not a
posterior distribution.

The accompanying code in \texttt{projects/density\_bioscales/src/}
drives table exports (\texttt{density\_summary.json},
\texttt{preset\_densities.csv}) and figures referenced in Section
``Results''. Buoyancy labels (\texttt{float}, \texttt{neutral},
\texttt{sink}) are qualitative helpers for comparing model body density
to a chosen still fluid. The intended use is methodological: keep
assumptions visible, keep outputs byte-reproducible, and separate
transparent fluid physics from deliberately coarse biomechanical
cartoons.
\end{frame}

\end{document}
